% TOP_PROBLEM: {"id": 20001890, "title": "Katok entropy rigidity and an exact time-change obstruction", "category_id": 6, "category": {"id": 6, "name": "geometry", "display_name": "Geometry", "description": "Euclidean and non-Euclidean geometry, geometric structures.", "slug": "geometry", "order_index": 6, "created_at": "2026-07-31T15:26:25.671Z"}, "status": "partially_solved"}
% FIRST_POSTED: 2026-09-27
% !TeX program = lualatex
% ATTEMPT_STATUS: unresolved
\documentclass[11pt]{article}
\usepackage[margin=25mm]{geometry}
\usepackage{fontspec}
\setmainfont{Latin Modern Roman}
\usepackage{amsmath,amssymb,mathtools}
\usepackage{unicode-math}
\setmathfont{Latin Modern Math}
\usepackage{xurl,hyperref}
\hypersetup{colorlinks=true,urlcolor=blue}
\setcounter{secnumdepth}{0}
\setlength{\parindent}{0pt}
\setlength{\parskip}{5pt}
\setlength{\emergencystretch}{3em}
\begin{document}
\section*{\texorpdfstring{Katok entropy rigidity and an exact time-change obstruction}{Katok entropy rigidity and an exact time-change obstruction}}\label{20001890}
\subsection{Definitions and mathematical statement}
Let $(M,g)$ be a closed negatively curved Riemannian manifold, and let $\varphi_t$ be its geodesic flow on the unit tangent bundle $SM$. The Liouville probability measure $\mu_L$ is the normalized invariant measure supplied by the Riemannian metric. Write $h_{\mu_L}(\varphi_1)$ for the metric entropy of the time-one map and $h_{\rm top}(\varphi_1)$ for its topological entropy.

Katok's selected rigidity assertion is
\[
 h_{\mu_L}(\varphi_1)=h_{\rm top}(\varphi_1)
 \quad\Longleftrightarrow\quad (M,g)\text{ is locally symmetric}.
\]
Local symmetry means that the curvature tensor is parallel, $\nabla\operatorname{Rm}=0$, equivalently the universal covering metric is locally a symmetric-space metric. Equality is measured for the actual Liouville measure, not merely for some invariant measure of maximal entropy.
\par\smallskip\textit{Formulation and concept references:} \href{https://arxiv.org/abs/2409.11197}{[S1]}.

\subsection{Short English statement}
Does Liouville measure have maximal geodesic-flow entropy exactly when a closed negatively curved metric is locally symmetric?
\subsection{Research attempt: a Riccati deficit and the limitation of scalar curvature}
\textbf{Status.} The general rigidity implication is not proved here. We derive an exact nonnegative deficit identity from the unstable Riccati equation, recover the negatively curved surface case using an explicitly identified external entropy inequality, and show why forcing that deficit to vanish would give the wrong target in higher dimensions.

\textbf{Source audit.} Theorem 1 of the inspected version arXiv:2409.11197v3 of [S1] concerns metrics in a sufficiently small $C^N$ neighborhood of a fixed real or complex hyperbolic metric, in dimension at least three. It yields homothety to that metric under entropy equality. Its neighborhood and background-metric hypotheses cannot be dropped. For the surface entropy comparison we also inspected A. Katok, \emph{Entropy and closed geodesics}, Theorem B and Section 3, at \url{https://akatok.s3.amazonaws.com/pub/AKETDS1982.pdf}. That classical theorem is an external input below, not a result proved by the elementary Riccati calculation.

For clarity take $M$ connected, smooth, and of dimension $d\geq2$. This is the usual connected setting of the rigidity question. All entropies are for unit-speed geodesic flow per unit time, and Liouville measure is normalized to be a probability. A disconnected variant has extra component bookkeeping and is not implicitly covered by the following surface argument.

\subsubsection{1. The geometric differential equation}
Along a unit-speed geodesic $\gamma$, the normal Jacobi equation is
\[
 J''+R_{\dot\gamma}J=0,
 \qquad R_v(w)=R(w,v)v,\quad w\perp v.
\]
Use the curvature convention in which $\langle R_vw,w\rangle$ is the sectional curvature when $v,w$ are orthonormal. Let $Y(t)$ be an unstable Jacobi tensor and define
\[
 U(t)=Y'(t)Y(t)^{-1}.
\]
The standard unstable construction for a closed negatively curved manifold gives a symmetric positive definite $U(v)$, continuous in $v\in SM$, differentiable along flow lines, and satisfying
\[
 XU+U^2+R_v=0.                                         \tag{1}
\]
Here $X$ is differentiation along the geodesic flow, with the normal spaces identified by parallel transport. The equation itself follows by differentiating $Y'Y^{-1}$ and using the Jacobi equation. Existence and regularity of the unstable tensor are standard geometric inputs; we do not assume the stable or unstable distributions are globally smooth in transverse directions.

Write $m=d-1$ and $H=\operatorname{tr}U$. Taking the trace gives
\[
 XH+\operatorname{tr}(U^2)+\operatorname{Ric}(v,v)=0.
                                                               \tag{2}
\]
Although $H$ need not be globally differentiable in every direction, it is continuously differentiable along the flow. Invariance of $\mu_L$ implies $\int XH\,d\mu_L=0$: integrate the difference quotient of $H\circ\varphi_t-H$, and use the continuous derivative along the compact phase space. Thus
\[
 \int_{SM}\operatorname{tr}(U^2)\,d\mu_L
 =-\int_{SM}\operatorname{Ric}(v,v)\,d\mu_L.              \tag{3}
\]

We use the entropy formula for smooth invariant volume in a smooth Anosov geodesic flow:
\[
 h_L:=h_{\mu_L}(\varphi_1)=\int_{SM}H\,d\mu_L.            \tag{4}
\]
This is the usual Pesin entropy formula expressed through unstable Jacobi fields. The connection with $H$ is visible directly from
\[
 \log|\det Y(t)|-\log|\det Y(0)|
 =\int_0^t H(\varphi_s v)\,ds.
\]
Replacing this Jacobi determinant by the unstable Jacobian in a fixed bundle metric changes its logarithm by a bounded endpoint correction, which disappears after division by $t$. The identification of metric entropy with the positive Lyapunov sum remains the named external entropy theorem, not an elementary consequence of determinant differentiation.

\subsubsection{2. An exact variance identity}
Decompose $U$ into its scalar and trace-free parts:
\[
 U=\frac{H}{m}I+U_0,\qquad\operatorname{tr}U_0=0.
\]
Symmetry gives the pointwise identity
\[
 \operatorname{tr}(U^2)=\frac{H^2}{m}+\|U_0\|_{\mathrm{HS}}^2.
\]
Combining this with (3)--(4) proves
\[
 \boxed{
 -m\int_{SM}\operatorname{Ric}(v,v)\,d\mu_L-h_L^2
 =\operatorname{Var}_{\mu_L}(H)
  +m\int_{SM}\|U_0\|_{\mathrm{HS}}^2\,d\mu_L.}
                                                               \tag{5}
\]
Both terms on the right are nonnegative. This is stronger than simply applying Cauchy--Schwarz twice, since it records both sources of loss separately.

At a fixed basepoint, rotational symmetry of the unit sphere gives the average $v_iv_j=\delta_{ij}/d$. The base marginal of normalized Liouville measure is normalized Riemannian volume. Therefore
\[
 \int_{SM}\operatorname{Ric}(v,v)\,d\mu_L
 =\frac1{d\operatorname{Vol}_g(M)}\int_M\operatorname{Scal}_g\,dV_g.
                                                               \tag{6}
\]
Equations (5)--(6) yield the explicit upper bound
\[
 h_L^2\leq-
 \frac{d-1}{d\operatorname{Vol}_g(M)}
 \int_M\operatorname{Scal}_g\,dV_g.                       \tag{7}
\]
No topological entropy has entered this derivation. In particular, entropy equality in the question has not yet forced equality in (7).

If equality does hold in (7), (5) shows that $H$ is constant and $U_0=0$ almost everywhere. Continuity and full support of $\mu_L$ make these statements true everywhere. Then $U=aI$ with a constant $a>0$, and (1) gives
\[
 R_v=-a^2 I\quad\text{on }v^\perp\quad\text{for every }v\in SM.
\]
Every two-plane has sectional curvature $-a^2$. Hence equality in the Ricci-average bound proves constant sectional curvature. This is a rigorous sufficient condition, but is potentially stronger than local symmetry.

\subsubsection{3. The surface case and its genuinely external ingredient}
Now let $d=2$, and write $A=\operatorname{Area}_g(M)$. The endomorphism $U$ is a scalar $u>0$. Equation (1) becomes
\[
 Xu+u^2+K\circ\pi=0.
\]
Gauss--Bonnet and (3) imply
\[
 \int_{SM}u^2\,d\mu_L=-\frac1A\int_M K\,dA
 =-\frac{2\pi\chi(M)}A=:q>0.                            \tag{8}
\]
Thus $h_L^2\leq q$, and more precisely
\[
 q-h_L^2=\int_{SM}(u-h_L)^2\,d\mu_L.                    \tag{9}
\]
The second side of the argument is Katok's surface entropy-area inequality:
\[
 h_{\mathrm{top}}(\varphi_1)^2\geq
 -\frac{2\pi\chi(M)}A=q.                                \tag{10}
\]
This is the external geometric theorem specified in the source audit. We do not derive (10) from (8); the Riccati integral controls Liouville entropy in the opposite direction and is not a substitute for (10).

If $h_L=h_{\mathrm{top}}$, (8)--(10) force $h_L^2=q$, so (9) gives $u\equiv h_L$. The scalar Riccati equation then gives $K\equiv-h_L^2$. This supplies the equality-case deduction with all analytic equalities exposed. Conversely, a surface of constant curvature $-a^2$ has unstable growth $a$ and topological entropy $a$, so the two entropies coincide. The latter topological entropy value follows from the standard identification with exponential volume growth in its hyperbolic universal cover, or from the classical constant-curvature computation in Katok's paper.

The proof explains a special feature of surfaces: Gauss--Bonnet turns the average curvature into a fixed area-topology quantity, and a separate entropy theorem places topological entropy on the other side of that same quantity. In dimension greater than two, the scalar-curvature integral is not fixed by volume and topology in this manner.

\subsubsection{4. A higher-dimensional stress test: complex hyperbolic space}
A proposed extension might try to prove entropy equality implies that the right side of (5) vanishes. This would be too strong. Consider complex hyperbolic space of complex dimension $q\geq2$, normalized so its holomorphic sectional curvature is $-4$, and take a closed quotient. Its real dimension is $d=2q$. For a unit vector $v$, the normal direction $Jv$ has curvature $-4$, while the $2q-2$ remaining normal directions have curvature $-1$.

In a parallel frame along a geodesic, the unstable Jacobi equations are $y''-4y=0$ in the first direction and $y''-y=0$ in the others. The unstable solutions grow as $e^{2t}$ and $e^t$, respectively. Consequently
\[
 \operatorname{spec}U=(2,\underbrace{1,\ldots,1}_{2q-2}),
 \quad H=2q,\quad\operatorname{tr}(U^2)=2q+2.             \tag{11}
\]
These are constant eigenvalues, yet $U$ is not scalar. The locally symmetric entropy computation gives
\[
 h_L=h_{\mathrm{top}}=2q,
\]
while the deficit in (5) equals
\[
 (2q-1)(2q+2)-(2q)^2=2q-2>0.                            \tag{12}
\]
There is no contradiction: $\operatorname{Var}(H)=0$, but the trace-free term is positive. The target includes these metrics. Thus (12) refutes the intermediate claim ``entropy equality forces equality in the Ricci-average bound,'' even on the intended equality examples. It does not refute the conjecture.

The calculation also shows that proving only constant $H$ would be different from proving $U=aI$. Locally symmetric spaces can have several fixed unstable expansion rates. Any successful rigidity argument must accommodate that spectrum rather than impose isotropic expansion by an inappropriate scalar equality case.

\subsubsection{5. Normalization checks and the remaining route}
If $\widetilde g=c^2g$ for a constant $c>0$, the map $v\mapsto v/c$ identifies the two unit tangent bundles and changes unit-speed geodesic time by $t\mapsto t/c$. Thus both entropies scale by $c^{-1}$, normalized Liouville probabilities correspond, and entropy equality is unchanged. The terms in (5) scale by $c^{-2}$: curvature scales by that factor, while $U$ and its trace scale by $c^{-1}$. Therefore neither a volume normalization nor a different curvature normalization can remove the defect in (12).

Another tempting shortcut uses the existence of a measure of maximal entropy. Negative curvature supplies such a measure, but the conjecture asks whether it is the particular smooth Liouville measure. Existence alone proves no equality for $\mu_L$. Similarly, integrating (2) against another invariant measure would change the averages and would not give (4) for that measure without the corresponding entropy hypotheses.

A possible next step would relate the equality of these two distinguished measures to unstable Jacobians on every closed geodesic, and then extract curvature information from those orbitwise data. Neither the scalar average (3) nor the equality of two entropy numbers by itself establishes that orbitwise conclusion. Even an appropriate cohomological statement for the unstable potential would still require a rigidity theorem converting it into parallel curvature. Those are substantive missing inputs here, rather than formal consequences of (5).

The accompanying diagnostic checks (12) for $2\leq q\leq10$, the decomposition behind (5) for explicit symmetric matrices, and its real-hyperbolic zero-deficit case. Such checks protect signs and multiplicities; they cannot infer geometric rigidity from finitely many matrices. The surface deduction is complete relative to the named standard entropy inputs. In higher dimensions, the exact deficit identity is proved, but the scalar method stops because its equality condition excludes legitimate complex hyperbolic equality examples. The full target remains unresolved by this attempt.

\subsection{Sources}
\begin{itemize}
\item[S1]T. Humbert, Katok's entropy conjecture near real and complex hyperbolic metrics, abstract, revised 2025.
\url{https://arxiv.org/abs/2409.11197}.
\textit{Formulation source.}
\item[S2] A. Katok, Entropy and closed geodesics, Ergodic Theory and Dynamical Systems 2 (1982), Theorem B and Section 3; author-hosted copy.
\url{https://akatok.s3.amazonaws.com/pub/AKETDS1982.pdf}.
\item[S3] T. Humbert, version 3 of [S1], revised 20 October 2025, Theorem 1 and Section 2.3.
\url{https://arxiv.org/html/2409.11197v3}.
\end{itemize}
\end{document}
